\chapter{The Kidneys}

\section*{Definition and Overview}
The kidneys are two bean-shaped organs, each about the size of a fist, located at the back of the abdomen on either side of the spine. Their main job is to filter the blood, removing waste products and excess fluid to form urine, while keeping the salts and other substances the body needs. The kidneys also regulate blood pressure, produce hormones that stimulate red blood cell production and keep bones strong, and help balance the body's acid and minerals. They filter around 180 litres of blood daily. This chapter explains how the kidneys work, how to keep them healthy, and the signs of kidney problems.

\section*{Epidemiology}
Kidney disease is common and often silent. Chronic kidney disease affects around one in ten Australians, and most people with early disease do not know they have it. The leading causes are diabetes and high blood pressure, and Aboriginal and Torres Strait Islander peoples are affected at higher rates. Around 25,000 Australians are treated for end-stage kidney disease with dialysis or a transplant. Kidney disease is among the leading causes of hospitalisation, and its costs grow with the ageing population, making kidney health a national priority.

\section*{Aetiology and Pathophysiology}
Each kidney contains about a million tiny filtering units called nephrons, each with a glomerulus, a knot of blood vessels that filters the blood, and a tubule that reclaims useful substances. Blood enters under high pressure, and the glomerulus filters water, salts, and waste into the tubule, while proteins and blood cells are held back. The tubules reabsorb water and nutrients, concentrate waste, and adjust the balance of sodium, potassium, acid, and minerals under the control of hormones. The kidneys produce renin to regulate blood pressure, erythropoietin to stimulate red blood cell production, and activate vitamin D for bone health. When nephrons are lost, through diabetes, high blood pressure, or other disease, the remaining nephrons compensate until the loss is advanced, which is why kidney disease is silent for so long.

\section*{Clinical Features}
\begin{itemize}
    \item The kidneys filter waste, balance fluids and minerals, and regulate blood pressure
    \item They produce hormones for red blood cells and bone health
    \item Urine is formed continuously and stored in the bladder
    \item Healthy kidneys work with no noticeable sensation
    \item Signs of kidney problems: swelling, foamy urine, blood in the urine, and fatigue
    \item High blood pressure is both a cause and a sign of kidney disease
    \item Most kidney disease causes no symptoms until advanced
    \item Kidney function can be measured with blood and urine tests
\end{itemize}

\section*{Differential Diagnosis}
Kidney symptoms and disease have many causes.
\begin{itemize}
    \item \textbf{Diabetes:} the leading cause of chronic kidney disease
    \item \textbf{High blood pressure:} the second leading cause
    \item \textbf{Glomerulonephritis:} inflammation of the filters
    \item \textbf{Polycystic kidney disease:} an inherited condition with cysts
    \item \textbf{Kidney stones and infection:} acute kidney problems
    \item \textbf{Obstruction:} stones, prostate disease, and tumours
    \item \textbf{Medicines and toxins:} damage from painkillers and other drugs
\end{itemize}

\section*{When to Seek Medical Care}
People with diabetes, high blood pressure, or a family history of kidney disease should have regular kidney checks, and anyone with swelling, foamy or bloody urine, or persistent fatigue should be assessed. Urgent care is needed for a sudden drop in urine, severe swelling, breathlessness, or confusion. Medicines that can harm the kidneys should be used with care.

\textbf{Call 000 (or local emergency services) for:}
\begin{itemize}
    \item Difficulty breathing or severe swelling
    \item Passing very little or no urine
    \item Confusion, drowsiness, or a seizure
    \item Severe flank pain with fever
    \item Chest pain or collapse
\end{itemize}

\section*{Diagnosis and Investigations}
\begin{itemize}
    \item \textbf{Blood tests:} creatinine, eGFR, and electrolytes
    \item \textbf{Urine tests:} protein, blood, and infection
    \item \textbf{Blood pressure measurement:} at every review
    \item \textbf{Ultrasound:} kidney size, structure, and obstruction
    \item \textbf{CT and other imaging:} for stones, tumours, and anatomy
    \item \textbf{Biopsy:} for glomerular disease
    \item \textbf{Regular monitoring:} tracking function over time
\end{itemize}

\section*{Management}
\begin{itemize}
    \item \textbf{Controlling the cause:} blood pressure, diabetes, and cholesterol
    \item \textbf{Kidney-protecting medicines:} ACE inhibitors and ARBs
    \item \textbf{Lifestyle:} healthy diet, low salt, weight, activity, and stopping smoking
    \item \textbf{Fluid management:} avoiding dehydration and, in advanced disease, restricting fluid
    \item \textbf{Treatment of complications:} anaemia, bone disease, and heart risk
    \item \textbf{Dialysis and transplant:} for end-stage kidney disease
    \item \textbf{Medicine safety:} reviewing drugs that harm the kidneys
    \item \textbf{Specialist care:} nephrology for advanced or progressive disease
\end{itemize}

\section*{Complications}
\begin{itemize}
    \item Chronic kidney disease progressing to kidney failure
    \item High blood pressure and heart disease
    \item Anaemia and fatigue
    \item Bone disease and fractures
    \item Fluid overload with swelling and breathlessness
    \item Electrolyte disturbances, including dangerous potassium levels
    \item Malnutrition and reduced quality of life
\end{itemize}

\section*{Prognosis and Outcomes}
Kidneys can be protected. With good control of blood pressure and diabetes, the progression of kidney disease is slowed or halted in most people, and many never reach kidney failure. Early detection through simple tests is the key, and lifestyle measures protect the kidneys at all ages. For the minority who develop end-stage disease, dialysis and transplant provide effective treatment, and kidney health is restored by early prevention.

\section*{Prevention and Self-Care}
\begin{itemize}
    \item Have regular kidney function tests if you are at risk
    \item Control blood pressure and blood sugar
    \item Eat a healthy diet and limit salt
    \item Maintain a healthy weight and stay active
    \item Do not smoke and limit alcohol
    \item Drink enough fluid to stay hydrated
    \item Use painkillers and other medicines with care
    \item Attend health checks and follow advice
    \item Seek early assessment for kidney symptoms
\end{itemize}

\section*{Sources and Further Reading}
The following reputable sources provide further information for healthcare students and patients seeking reliable, current advice about the kidneys.
\begin{itemize}
    \item Kidney Health Australia. \textit{How your kidneys work.} https://kidney.org.au/
    \item Healthdirect Australia. \textit{Kidneys and kidney health.} https://www.healthdirect.gov.au/
    \item Better Health Channel. \textit{Kidneys.} https://www.betterhealth.vic.gov.au/
    \item Diabetes Australia. \textit{Kidney health and diabetes.} https://www.diabetesaustralia.com.au/
    \item NHS. \textit{Chronic kidney disease.} https://www.nhs.uk/
    \item National Kidney Foundation. \textit{Kidney education.} https://www.kidney.org/
\end{itemize}
